\documentclass[10pt,a4paper]{amsart}
\usepackage{amsmath,amssymb,amsthm,graphicx}
\usepackage{verbatim}
\newenvironment{reference}{\comment}{\endcomment}
%\newenvironment{reference}{}{}
\newenvironment{slogan}{\comment}{\endcomment}
\newenvironment{history}{\comment}{\endcomment}

% For commutative diagrams we use Xy-pic
\usepackage[all]{xy}

% We use 2cell for 2-commutative diagrams.
\xyoption{2cell}
% Theorem environments.
%
\theoremstyle{plain}
\newtheorem{theorem}[subsection]{Theorem}
\newtheorem{proposition}[subsection]{Proposition}
\newtheorem{lemma}[subsection]{Lemma}

\theoremstyle{definition}
\newtheorem{definition}[subsection]{Definition}
\newtheorem{example}[subsection]{Example}
\newtheorem{exercise}[subsection]{Exercise}
\newtheorem{situation}[subsection]{Situation}

\theoremstyle{remark}
\newtheorem{remark}[subsection]{Remark}
\newtheorem{remarks}[subsection]{Remarks}

\numberwithin{equation}{subsection}

% Macros
%
\def\lim{\mathop{\mathrm{lim}}\nolimits}
\def\colim{\mathop{\mathrm{colim}}\nolimits}
\def\Spec{\mathop{\mathrm{Spec}}}
\def\Hom{\mathop{\mathrm{Hom}}\nolimits}
\def\Ext{\mathop{\mathrm{Ext}}\nolimits}
\def\SheafHom{\mathop{\mathcal{H}\!\mathit{om}}\nolimits}
\def\SheafExt{\mathop{\mathcal{E}\!\mathit{xt}}\nolimits}
\def\Sch{\mathit{Sch}}
\def\Mor{\mathop{\mathrm{Mor}}\nolimits}
\def\Ob{\mathop{\mathrm{Ob}}\nolimits}
\def\Sh{\mathop{\mathit{Sh}}\nolimits}
\def\NL{\mathop{N\!L}\nolimits}
\def\CH{\mathop{\mathrm{CH}}\nolimits}
\def\proetale{{pro\text{-}\acute{e}tale}}
\def\etale{{\acute{e}tale}}
\def\QCoh{\mathit{QCoh}}
\def\Ker{\mathop{\mathrm{Ker}}}
\def\Im{\mathop{\mathrm{Im}}}
\def\Coker{\mathop{\mathrm{Coker}}}
\def\Coim{\mathop{\mathrm{Coim}}}

% Boxtimes
%
\DeclareMathSymbol{\boxtimes}{\mathbin}{AMSa}{"02}

%
% Macros for moduli stacks/spaces
%
\def\QCohstack{\mathcal{QC}\!\mathit{oh}}
\def\Cohstack{\mathcal{C}\!\mathit{oh}}
\def\Spacesstack{\mathcal{S}\!\mathit{paces}}
\def\Quotfunctor{\mathrm{Quot}}
\def\Hilbfunctor{\mathrm{Hilb}}
\def\Curvesstack{\mathcal{C}\!\mathit{urves}}
\def\Polarizedstack{\mathcal{P}\!\mathit{olarized}}
\def\Complexesstack{\mathcal{C}\!\mathit{omplexes}}
% \Pic is the operator that assigns to X its picard group, usage \Pic(X)
% \Picardstack_{X/B} denotes the Picard stack of X over B
% \Picardfunctor_{X/B} denotes the Picard functor of X over B
\def\Pic{\mathop{\mathrm{Pic}}\nolimits}
\def\Picardstack{\mathcal{P}\!\mathit{ic}}
\def\Picardfunctor{\mathrm{Pic}}
\def\Deformationcategory{\mathcal{D}\!\mathit{ef}}
\usepackage{cleveref}
\crefname{originaldefinition}{original definition}{original definitions}
\Crefname{originaldefinition}{Original definition}{Original definitions}
\makeatletter
\newlabel{formal-spaces-definition-formal-algebraic-space}{{0AIM}{source}{}{sourcepointer-1}{}}
\newlabel{formal-spaces-definition-formal-algebraic-space@cref}{{[originaldefinition][0][]0AIM}{[1][0][]source}{}{}{}}
\makeatother
\crefname{originalsection}{original section}{original sections}
\Crefname{originalsection}{Original section}{Original sections}
\makeatletter
\newlabel{formal-spaces-example-reduction-affine-formal-spectrum}{{0GB6}{source}{}{sourcepointer-2}{}}
\newlabel{formal-spaces-example-reduction-affine-formal-spectrum@cref}{{[originalsection][0][]0GB6}{[1][0][]source}{}{}{}}
\makeatother
\crefname{originallemma}{original lemma}{original lemmas}
\Crefname{originallemma}{Original lemma}{Original lemmas}
\makeatletter
\newlabel{formal-spaces-lemma-property-defines-property-morphisms}{{0ANG}{source}{}{sourcepointer-3}{}}
\newlabel{formal-spaces-lemma-property-defines-property-morphisms@cref}{{[originallemma][0][]0ANG}{[1][0][]source}{}{}{}}
\makeatother
\crefname{originallemma}{original lemma}{original lemmas}
\Crefname{originallemma}{Original lemma}{Original lemmas}
\makeatletter
\newlabel{formal-spaces-lemma-reduction-formal-algebraic-space}{{0AIN}{source}{}{sourcepointer-4}{}}
\newlabel{formal-spaces-lemma-reduction-formal-algebraic-space@cref}{{[originallemma][0][]0AIN}{[1][0][]source}{}{}{}}
\makeatother
\crefname{originallemma}{original lemma}{original lemmas}
\Crefname{originallemma}{Original lemma}{Original lemmas}
\makeatletter
\newlabel{formal-spaces-lemma-reduction-smooth}{{0GB7}{source}{}{sourcepointer-5}{}}
\newlabel{formal-spaces-lemma-reduction-smooth@cref}{{[originallemma][0][]0GB7}{[1][0][]source}{}{}{}}
\makeatother
\crefname{originalreference}{original reference}{original references}
\Crefname{originalreference}{Original reference}{Original references}
\makeatletter
\newlabel{lemma-finite-type-red}{{0GC2}{source}{}{sourcepointer-6}{}}
\newlabel{lemma-finite-type-red@cref}{{[originalreference][0][]0GC2}{[1][0][]source}{}{}{}}
\makeatother
\crefname{originalsection}{original section}{original sections}
\Crefname{originalsection}{Original section}{Original sections}
\makeatletter
\newlabel{spaces-morphisms-section-finite-type}{{03XE}{source}{}{sourcepointer-7}{}}
\newlabel{spaces-morphisms-section-finite-type@cref}{{[originalsection][0][]03XE}{[1][0][]source}{}{}{}}
\makeatother
\title{Finite type of a morphism of formal spaces is finite type of the reductions}
\author{The Stacks Project Authors}
\date{stacks-project:restricted.tex:local-snapshot-20261012; GNU FDL 1.2 or later}
\begin{document}
\maketitle
\noindent\textit{Editorial note.} All delivered mathematical text is the author's own, copied from the pinned Stacks Project chapter file restricted.tex; no statement or proof step was rewritten, reordered or removed.\par\medskip
\noindent\textit{Editorial note.} The counted proof is the author's complete argument: the proof derives the (1)-(2)-(3) equivalence from the property-transfer lemma, computes the reduction of the affine pieces as the quotient by the ideals of topologically nilpotent elements, and closes via the etale-local nature of local finite type\par\medskip
\noindent\textit{Editorial note.} The label lemma-finite-type-red-morphisms names the counted statement, and the proof environment immediately following it in the source is delivered verbatim as its proof.\par\medskip
\noindent\textit{Editorial note.} The counted statement is carried by the author's own lemma environment in the chapter source.\par\medskip
\noindent\textit{Allowed standard prerequisites (editorial).}\begin{itemize}
\item the ring-theoretic property of arrows of WAdm\textasciicircum{}count defined in lemma-finite-type-red of the same chapter
\item the correspondence between properties of arrows of WAdm\textasciicircum{}count and properties of morphisms of formal algebraic spaces (Formal Spaces, Lemma formal-spaces-lemma-property-defines-property-morphisms)
\item the reduction functor on formal algebraic spaces and its etale-local behaviour (Formal Spaces, Lemmas formal-spaces-lemma-reduction-formal-algebraic-space and formal-spaces-lemma-reduction-smooth, Example formal-spaces-example-reduction-affine-formal-spectrum)
\item local finite type of morphisms of algebraic spaces is etale local on source and target (Morphisms of Spaces, Section spaces-morphisms-section-finite-type)
\end{itemize}\medskip
\ifdefined\hypertarget\hypertarget{sourcepointer-1}{}\fi
\noindent\textit{Source reference.} Original definition 0AIM: sources/stacks-chow/tags.txt, source lines 13270--13270. Prerequisite pointer invoked at proof line(s) 2813; it names a standard fact of a sibling chapter that the counted argument uses, and every step of the counted proof built on it is retained verbatim in the delivered body.\par\medskip
\ifdefined\hypertarget\hypertarget{sourcepointer-2}{}\fi
\noindent\textit{Source reference.} Original section 0GB6: sources/stacks-chow/tags.txt, source lines 20361--20361. Prerequisite pointer invoked at proof line(s) 2854; it names a standard fact of a sibling chapter that the counted argument uses, and every step of the counted proof built on it is retained verbatim in the delivered body.\par\medskip
\ifdefined\hypertarget\hypertarget{sourcepointer-3}{}\fi
\noindent\textit{Source reference.} Original lemma 0ANG: sources/stacks-chow/tags.txt, source lines 13439--13439. Prerequisite pointer invoked at proof line(s) 2837; it names a standard fact of a sibling chapter that the counted argument uses, and every step of the counted proof built on it is retained verbatim in the delivered body.\par\medskip
\ifdefined\hypertarget\hypertarget{sourcepointer-4}{}\fi
\noindent\textit{Source reference.} Original lemma 0AIN: sources/stacks-chow/tags.txt, source lines 13271--13271. Prerequisite pointer invoked at proof line(s) 2844; it names a standard fact of a sibling chapter that the counted argument uses, and every step of the counted proof built on it is retained verbatim in the delivered body.\par\medskip
\ifdefined\hypertarget\hypertarget{sourcepointer-5}{}\fi
\noindent\textit{Source reference.} Original lemma 0GB7: sources/stacks-chow/tags.txt, source lines 20362--20362. Prerequisite pointer invoked at proof line(s) 2846; it names a standard fact of a sibling chapter that the counted argument uses, and every step of the counted proof built on it is retained verbatim in the delivered body.\par\medskip
\ifdefined\hypertarget\hypertarget{sourcepointer-6}{}\fi
\noindent\textit{Source reference.} Original reference 0GC2: sources/stacks-restricted/pinned-original-source.tex, source lines 2631--2631. Prerequisite pointer invoked at proof line(s) 2810; it names a standard fact of the same chapter that the counted argument uses, and every step of the counted proof built on it is retained verbatim in the delivered body.\par\medskip
\ifdefined\hypertarget\hypertarget{sourcepointer-7}{}\fi
\noindent\textit{Source reference.} Original section 03XE: sources/stacks-chow/tags.txt, source lines 5145--5145. Prerequisite pointer invoked at proof line(s) 2862; it names a standard fact of a sibling chapter that the counted argument uses, and every step of the counted proof built on it is retained verbatim in the delivered body.\par\medskip
\begin{lemma}
\label{lemma-finite-type-red-morphisms}
Let $S$ be a scheme. Let $f : X \to Y$ be a morphism of
locally countably indexed formal algebraic spaces over $S$.
The following are equivalent
\begin{enumerate}
\item for every commutative diagram
$$
\xymatrix{
U \ar[d] \ar[r] & V \ar[d] \\
X \ar[r] & Y
}
$$
with $U$ and $V$ affine formal algebraic spaces, $U \to X$ and $V \to Y$
representable by algebraic spaces and \'etale, the morphism $U \to V$
corresponds to an arrow of $\textit{WAdm}^{count}$ satisfying the
property defined in Lemma \ref{lemma-finite-type-red},
\item there exists a covering $\{Y_j \to Y\}$ as in
Formal Spaces,
Definition \ref{formal-spaces-definition-formal-algebraic-space}
and for each $j$
a covering $\{X_{ji} \to Y_j \times_Y X\}$ as in
Formal Spaces,
Definition \ref{formal-spaces-definition-formal-algebraic-space}
such that each $X_{ji} \to Y_j$  corresponds
to an arrow of $\textit{WAdm}^{count}$ satisfying the
property defined in Lemma \ref{lemma-finite-type-red},
\item there exist a covering $\{X_i \to X\}$ as in
Formal Spaces,
Definition \ref{formal-spaces-definition-formal-algebraic-space}
and for each $i$ a factorization $X_i \to Y_i \to Y$ where $Y_i$
is an affine formal algebraic space, $Y_i \to Y$ is representable
by algebraic spaces and \'etale, and $X_i \to Y_i$ corresponds
to an arrow of $\textit{WAdm}^{count}$ satisfying the
property defined in Lemma \ref{lemma-finite-type-red}, and
\item the morphism $f_{red} : X_{red} \to Y_{red}$ is locally of finite type.
\end{enumerate}
\end{lemma}

\begin{proof}
The equivalence of (1), (2), and (3) follows from
Lemma \ref{lemma-finite-type-red} and an application of
Formal Spaces, Lemma
\ref{formal-spaces-lemma-property-defines-property-morphisms}.
Let $Y_j$ and $X_{ji}$ be as in (2). Then
\begin{itemize}
\item The families $\{Y_{j, red} \to Y_{red}\}$ and
$\{X_{ji, red} \to X_{red}\}$ are \'etale coverings by affine schemes.
This follows from the discussion in the proof of
Formal Spaces, Lemma
\ref{formal-spaces-lemma-reduction-formal-algebraic-space}
or directly from
Formal Spaces, Lemma \ref{formal-spaces-lemma-reduction-smooth}.
\item If $X_{ji} \to Y_j$ corresponds to the morphism
$B_j \to A_{ji}$ of $\textit{WAdm}^{count}$, then
$X_{ji, red} \to Y_{j, red}$ corresponds to the ring map
$B_j/\mathfrak b_j \to A_{ji}/\mathfrak a_{ji}$ where
$\mathfrak b_j$ and $\mathfrak a_{ji}$ are the ideals of
topologically nilpotent elements. This follows from
Formal Spaces, Example
\ref{formal-spaces-example-reduction-affine-formal-spectrum}.
Hence $X_{ji, red} \to Y_{j, red}$ is locally of finite type
if and only if $B_j \to A_{ji}$ satisfies the
property defined in Lemma \ref{lemma-finite-type-red}.
\end{itemize}
The equivalence of (2) and (4) follows from these remarks because
being locally of finite type is a property of morphisms of algebraic
spaces which is \'etale local on source and target, see discussion in
Morphisms of Spaces, Section \ref{spaces-morphisms-section-finite-type}.
\end{proof}

\end{document}
